% Options for packages loaded elsewhere
\PassOptionsToPackage{unicode}{hyperref}
\PassOptionsToPackage{hyphens}{url}
%
\documentclass[
  doc,floatsintext]{apa6}
\usepackage{amsmath,amssymb}
\usepackage{iftex}
\ifPDFTeX
  \usepackage[T1]{fontenc}
  \usepackage[utf8]{inputenc}
  \usepackage{textcomp} % provide euro and other symbols
\else % if luatex or xetex
  \usepackage{unicode-math} % this also loads fontspec
  \defaultfontfeatures{Scale=MatchLowercase}
  \defaultfontfeatures[\rmfamily]{Ligatures=TeX,Scale=1}
\fi
\usepackage{lmodern}
\ifPDFTeX\else
  % xetex/luatex font selection
\fi
% Use upquote if available, for straight quotes in verbatim environments
\IfFileExists{upquote.sty}{\usepackage{upquote}}{}
\IfFileExists{microtype.sty}{% use microtype if available
  \usepackage[]{microtype}
  \UseMicrotypeSet[protrusion]{basicmath} % disable protrusion for tt fonts
}{}
\makeatletter
\@ifundefined{KOMAClassName}{% if non-KOMA class
  \IfFileExists{parskip.sty}{%
    \usepackage{parskip}
  }{% else
    \setlength{\parindent}{0pt}
    \setlength{\parskip}{6pt plus 2pt minus 1pt}}
}{% if KOMA class
  \KOMAoptions{parskip=half}}
\makeatother
\usepackage{xcolor}
\usepackage{graphicx}
\makeatletter
\def\maxwidth{\ifdim\Gin@nat@width>\linewidth\linewidth\else\Gin@nat@width\fi}
\def\maxheight{\ifdim\Gin@nat@height>\textheight\textheight\else\Gin@nat@height\fi}
\makeatother
% Scale images if necessary, so that they will not overflow the page
% margins by default, and it is still possible to overwrite the defaults
% using explicit options in \includegraphics[width, height, ...]{}
\setkeys{Gin}{width=\maxwidth,height=\maxheight,keepaspectratio}
% Set default figure placement to htbp
\makeatletter
\def\fps@figure{htbp}
\makeatother
\setlength{\emergencystretch}{3em} % prevent overfull lines
\providecommand{\tightlist}{%
  \setlength{\itemsep}{0pt}\setlength{\parskip}{0pt}}
\setcounter{secnumdepth}{-\maxdimen} % remove section numbering
% Make \paragraph and \subparagraph free-standing
\makeatletter
\ifx\paragraph\undefined\else
  \let\oldparagraph\paragraph
  \renewcommand{\paragraph}{
    \@ifstar
      \xxxParagraphStar
      \xxxParagraphNoStar
  }
  \newcommand{\xxxParagraphStar}[1]{\oldparagraph*{#1}\mbox{}}
  \newcommand{\xxxParagraphNoStar}[1]{\oldparagraph{#1}\mbox{}}
\fi
\ifx\subparagraph\undefined\else
  \let\oldsubparagraph\subparagraph
  \renewcommand{\subparagraph}{
    \@ifstar
      \xxxSubParagraphStar
      \xxxSubParagraphNoStar
  }
  \newcommand{\xxxSubParagraphStar}[1]{\oldsubparagraph*{#1}\mbox{}}
  \newcommand{\xxxSubParagraphNoStar}[1]{\oldsubparagraph{#1}\mbox{}}
\fi
\makeatother
% definitions for citeproc citations
\NewDocumentCommand\citeproctext{}{}
\NewDocumentCommand\citeproc{mm}{%
  \begingroup\def\citeproctext{#2}\cite{#1}\endgroup}
\makeatletter
 % allow citations to break across lines
 \let\@cite@ofmt\@firstofone
 % avoid brackets around text for \cite:
 \def\@biblabel#1{}
 \def\@cite#1#2{{#1\if@tempswa , #2\fi}}
\makeatother
\newlength{\cslhangindent}
\setlength{\cslhangindent}{1.5em}
\newlength{\csllabelwidth}
\setlength{\csllabelwidth}{3em}
\newenvironment{CSLReferences}[2] % #1 hanging-indent, #2 entry-spacing
 {\begin{list}{}{%
  \setlength{\itemindent}{0pt}
  \setlength{\leftmargin}{0pt}
  \setlength{\parsep}{0pt}
  % turn on hanging indent if param 1 is 1
  \ifodd #1
   \setlength{\leftmargin}{\cslhangindent}
   \setlength{\itemindent}{-1\cslhangindent}
  \fi
  % set entry spacing
  \setlength{\itemsep}{#2\baselineskip}}}
 {\end{list}}
\usepackage{calc}
\newcommand{\CSLBlock}[1]{\hfill\break\parbox[t]{\linewidth}{\strut\ignorespaces#1\strut}}
\newcommand{\CSLLeftMargin}[1]{\parbox[t]{\csllabelwidth}{\strut#1\strut}}
\newcommand{\CSLRightInline}[1]{\parbox[t]{\linewidth - \csllabelwidth}{\strut#1\strut}}
\newcommand{\CSLIndent}[1]{\hspace{\cslhangindent}#1}
\ifLuaTeX
\usepackage[bidi=basic]{babel}
\else
\usepackage[bidi=default]{babel}
\fi
\babelprovide[main,import]{english}
% get rid of language-specific shorthands (see #6817):
\let\LanguageShortHands\languageshorthands
\def\languageshorthands#1{}
% Manuscript styling
\usepackage{upgreek}
\captionsetup{font=singlespacing,justification=justified}

% Table formatting
\usepackage{longtable}
\usepackage{lscape}
% \usepackage[counterclockwise]{rotating}   % Landscape page setup for large tables
\usepackage{multirow}		% Table styling
\usepackage{tabularx}		% Control Column width
\usepackage[flushleft]{threeparttable}	% Allows for three part tables with a specified notes section
\usepackage{threeparttablex}            % Lets threeparttable work with longtable

% Create new environments so endfloat can handle them
% \newenvironment{ltable}
%   {\begin{landscape}\centering\begin{threeparttable}}
%   {\end{threeparttable}\end{landscape}}
\newenvironment{lltable}{\begin{landscape}\centering\begin{ThreePartTable}}{\end{ThreePartTable}\end{landscape}}

% Enables adjusting longtable caption width to table width
% Solution found at http://golatex.de/longtable-mit-caption-so-breit-wie-die-tabelle-t15767.html
\makeatletter
\newcommand\LastLTentrywidth{1em}
\newlength\longtablewidth
\setlength{\longtablewidth}{1in}
\newcommand{\getlongtablewidth}{\begingroup \ifcsname LT@\roman{LT@tables}\endcsname \global\longtablewidth=0pt \renewcommand{\LT@entry}[2]{\global\advance\longtablewidth by ##2\relax\gdef\LastLTentrywidth{##2}}\@nameuse{LT@\roman{LT@tables}} \fi \endgroup}

% \setlength{\parindent}{0.5in}
% \setlength{\parskip}{0pt plus 0pt minus 0pt}

% Overwrite redefinition of paragraph and subparagraph by the default LaTeX template
% See https://github.com/crsh/papaja/issues/292
\makeatletter
\renewcommand{\paragraph}{\@startsection{paragraph}{4}{\parindent}%
  {0\baselineskip \@plus 0.2ex \@minus 0.2ex}%
  {-1em}%
  {\normalfont\normalsize\bfseries\itshape\typesectitle}}

\renewcommand{\subparagraph}[1]{\@startsection{subparagraph}{5}{1em}%
  {0\baselineskip \@plus 0.2ex \@minus 0.2ex}%
  {-\z@\relax}%
  {\normalfont\normalsize\itshape\hspace{\parindent}{#1}\textit{\addperi}}{\relax}}
\makeatother

\makeatletter
\usepackage{etoolbox}
\patchcmd{\maketitle}
  {\section{\normalfont\normalsize\abstractname}}
  {\section*{\normalfont\normalsize\abstractname}}
  {}{\typeout{Failed to patch abstract.}}
\patchcmd{\maketitle}
  {\section{\protect\normalfont{\@title}}}
  {\section*{\protect\normalfont{\@title}}}
  {}{\typeout{Failed to patch title.}}
\makeatother

\usepackage{xpatch}
\makeatletter
\xapptocmd\appendix
  {\xapptocmd\section
    {\addcontentsline{toc}{section}{\appendixname\ifoneappendix\else~\theappendix\fi: #1}}
    {}{\InnerPatchFailed}%
  }
{}{\PatchFailed}
\makeatother
\keywords{Trust in science; science knowledge; deficit model}
\usepackage{csquotes}
\usepackage{placeins} 

\ifLuaTeX
  \usepackage{selnolig}  % disable illegal ligatures
\fi
\usepackage{bookmark}
\IfFileExists{xurl.sty}{\usepackage{xurl}}{} % add URL line breaks if available
\urlstyle{same}
\hypersetup{
  pdftitle={The rational impression account of trust in science},
  pdfauthor={Jan Pfänder1, Sophie de Rouilhan1, \& Hugo Mercier1},
  pdflang={en-EN},
  pdfkeywords={Trust in science; science knowledge; deficit model},
  hidelinks,
  pdfcreator={LaTeX via pandoc}}

\title{The rational impression account of trust in science}
\author{Jan Pfänder\textsuperscript{1}, Sophie de Rouilhan\textsuperscript{1}, \& Hugo Mercier\textsuperscript{1}}
\date{}


\shorttitle{Rational impression account}

\authornote{

Correspondence concerning this article should be addressed to Hugo Mercier, . E-mail: \href{mailto:hugo.mercier@gmail.com}{\nolinkurl{hugo.mercier@gmail.com}}

}

\affiliation{\vspace{0.5cm}\textsuperscript{1} Institut Jean Nicod, Département d'études cognitives, ENS, EHESS, PSL University, CNRS, France}

\abstract{%
Trust in science matters to a range of significant outcomes such as intent to vaccinate and belief in climate change. Around the globe, most people trust science at least to some extent. However, the causes of this trust aren't well understood. Here we propose a `rational impression account' of trust in science: people trust scientists because they are impressed by their findings. This impression of trust persists even after knowledge of the specific content has vanished. We present evidence for this account in two experiments (total n = 696) with UK participants. In Experiment 1, we find that impressive scientific findings lead participants to think of the scientists as more competent and their scientific discipline as more trustworthy. In Experiment 2, we show that participants have these impressions despite forgetting the content that generated them. The rational impression account can explain why people trust science without understanding much of it. It also stresses the relevance of science communication.
}



\begin{document}
\maketitle

\section{Introduction}\label{introduction}

Jon D. Miller, a pioneer in measuring science literacy, found that at the end of the 20th century, approximately 17\% of Americans qualified as scientifically literate--a measure he described as the ``level of skill required to read most of the articles in the Tuesday science section of The New York Times'' (Miller, 2004). His estimate for the early 1990s had been even lower, around 10\%. By contrast, the level of trust in science in the US has been relatively high and remarkably stable (Funk \& Kennedy, 2020; Funk, Tyson, Kennedy, \& Johnson, 2020; Smith \& Son, 2013): Since the 1970s, on average 40\% of Americans say they have a great deal of confidence in the scientific community, the second highest score among 13 institutions included in the US General Social Survey (GSS), just behind the military. How can we explain this gap between understanding science and trusting it? Here, we argue that people might trust science in large part because they have been impressed by it.

We begin by presenting global evidence that most people trust science at least to some extent. We then review existing explanations for why people trust science, in particular the deficit-model, as well as empirical observations that seem to violate their predictions. We proceed by proposing an alternative model, which we term the `rational impression account'. In two experiments, we provide evidence in favor of this account.

Trust in science matters. It is related to many desirable outcomes, from acceptance of anthropogenic climate change (Cologna \& Siegrist, 2020) or vaccination (Lindholt, Jørgensen, Bor, \& Petersen, 2021; Sturgis, Brunton-Smith, \& Jackson, 2021) to following recommendations during COVID (Algan, Cohen, Davoine, Foucault, \& Stantcheva, 2021).

Across the globe, most people report trusting science at least to some extent (Wellcome Global Monitor, 2018, 2020). Recently, a large survey project in 68 countries found trust in scientists to be ``moderately high'' across all countries (mean = 3.62; sd= 0.70; Scale: 1 = very low, 2 = somewhat low, 3 = neither high nor low, 4 = somewhat high, 5 = very high), with not a single country below mid-point trust. A recent study in the US found that participants almost always accepted the scientific consensus on basic, non-politicized knowledge questions (e.g.~`Are electrons are smaller than atoms?'), even those participants who said they do not trust science and who believed anti-science conspiracy theories, e.g.~that the earth is flat (Pfänder, Kerzreho, \& Mercier, 2024).

How can we explain trust in science? The literature on the public understanding of science has generally focused not on explaining trust, but a lack of it, and more generally negative attitudes towards science (Bauer, Allum, \& Miller, 2007).

One explanation is what Gauchat (2011) coined the \emph{alienation model}. According to the alienation model, the ``public disassociation with science is a symptom of a general disenchantment with late modernity, mainly, the limitations associated with codified expertise, rational bureaucracy, and institutional authority'' (Gauchat, 2011). This explanation builds on the work of social theorists (Beck, 1992; Giddens, 1991; Habermas, 1989; see Gauchat, 2011 for an overview). According to these theorists, a modern, complex world increasingly requires expertise, and thus shapes institutions of knowledge elites. People who are not part of these institutions experience a lack of agency, resulting in a feeling of alienation. If the alienation model might account for some of the distrust towards science, its goal is not to explain why most people still trust science.

A second explanation which has dominated the literature on public understanding of science for decades (Bauer et al., 2007), is widely known as the ``deficit model''. According to the deficit model, people do not trust science enough, because they do not know enough about it. The deficit model implies that when people trust science, it is because they know a lot about it (Bauer et al., 2007).

Seemingly in line with that view, large scale survey studies consistently identify education, and in particular science education, to be the strongest correlate of trust in science (Bak, 2001; Noy \& O'Brien, 2019; Wellcome Global Monitor, 2018, 2020): More educated people tend to trust science more. The deficit model provides a straightforward causal explanation for this correlation: education fosters trust in science by transmitting knowledge and improving understanding of science.

Evidence for this causal model, however, is far from conclusive. While there is no consensus among researchers on how to measure understanding of science (Gauchat \& Andrews, 2018), a widely used measure consists of asking people a set of basic science knowledge questions (see e.g. Durant, Evans, \& Thomas, 1989; Miller, 1998). Past research has found science knowledge to be alarming low (Miller, 2004; National Academies of Sciences, Engineering, and Medicine, 2016), contrasting with the relatively high levels of trust in science. Moreover, science knowledge is only weakly associated with attitudes towards science (Allum, Sturgis, Tabourazi, \& Brunton-Smith, 2008). Other research has shown that the correlation between education and positive attitudes towards science holds even after controlling for science knowledge, which led the researchers to conclude that only part of the effect of education could be explained by transmitting science knowledge (Bak, 2001).

\section{The rational impression account}\label{the-rational-impression-account}

Here, we seek to explain the following stylized facts: (i) Globally, most people have at least some trust in science, and almost everyone, at least in the US, trusts (nearly all of) basic science; (ii) Education, and in particular science education, is the strongest predictor of trust in science; (iii) The levels of science knowledge are relatively low, and the association between science knowledge and trust in science is weak.

All of these observations are compatible with a model of trust in science which we dub the `rational impression account'. According to this account, people trust science because they are impressed by it, but they forget the details of the findings that impressed them.

For an information to be impressive, at least two criteria should be met: (i) it is perceived as hard to uncover and (ii) there is reason to believe its true. By (i), we mean information that most people would have no idea how to find out themselves. For example, most people would probably not be very impressed by someone telling them the biggest local park tree has exactly 110,201 leaves. Even though obtaining this information implies an exhausting counting effort, everyone in principle knows how to do it. By contrast, finding out that it takes light \href{https://imagine.gsfc.nasa.gov/features/cosmic/milkyway_info.html}{approximately 100'000 years to travel from one end of the Milky Way to the other} is probably impressive to most people, as they would not know how such a distance can be measured. Regarding (ii), people can sometimes judge the accuracy of scientific findings for themselves, for example when they are exposed to accessible and convincing explanations in school, e.g.~simple explanations that have a broad explanatory scope (T. Lombrozo, 2007; Read \& Marcus-Newhall, 1993; for a review, see Tania Lombrozo, 2006). When people cannot evaluate the quality of the information for themselves, they need to rely on other cues to establish whether it might be true or not. One such cue is the degree of consensus between different sources. It has been shown that when they lack relevant background knowledge, people rely on social consensus as a heuristic to judge the accuracy of a piece of information and the competence of its source (Pfänder, De Courson, \& Mercier, 2025). This heuristic is rational, in the sense that it leads to sound inferences if the sources are independent and unbiased. As a result, unless we suspect a conspiracy among scientists, we should be impressed when they agree on something like the length of the Milky Way, and think of the scientists as competent, even if we have no idea how they got to agree on this measure.

Once people are impressed, that impression tends to stick, even if specific content is forgotten. We commonly form impressions of the people around us while forgetting the details of how we formed these impressions: If a colleague fixes our computer, we might forget exactly how they fixed it, yet remember that they are good at fixing computers. For an extreme example, it has been demonstrated that patients with severe amnesia continued to experience emotions linked to events they could not recall (Feinstein, Duff, \& Tranel, 2010). More relevantly here, a study has shown that while people find some science-related explanations more satisfying than others, this did not predict how well they could recall the explanations shortly after (Liquin \& Lombrozo, 2022), suggesting that that impressions and knowledge can be quite detached.

The rational impression account of trust in science can reconcile the stylized facts outlined above: For (i), most people in the world receive at least some basic education that exposes them to impressive scientific findings, setting up a baseline of trust in science. For (ii), education can be assumed to be the primary place of exposure to science, making it the most important predictor. For (iii), specific knowledge about or understanding of science is forgotten, while the impression of competence and trust is maintained.

While the rational impression account, we argue, can reconcile well-established facts about trust in science and science knowledge, there is no experimental evidence directly testing its central mechanism. This is the goal of the present studies.

\section{The present studies}\label{the-present-studies}

We tested the predictions of the rational impression account in two experiments (total n = 696): In Experiment 1, we tested whether impressive scientific findings lead people to think of the scientists as more competent and trustworthy. In Experiment 2, we tested whether that impression persists even if specific content is forgotten.

Both experiments were preregistered, and the choice of sample size was informed by power simulations. All materials, data and code can be found on Open Science Framework \href{https://osf.io/j3bk4/}{project page} (\url{https://osf.io/j3bk4/}). All analyses were conducted in R (version 4.2.2) using R Studio. We used two-sided tests for all hypotheses. Unless mentioned otherwise, we report unstandardized estimates that can be interpreted in units of the scales we used to measure the variables.

As part of this project, we conducted two additional experiments which we present in detail in the ESM. The first experiment (which we refer to as `Experiment 1b' in the ESM) is almost identical to Experiment 1, but suffered from a minor technical error during the implementation. Nonetheless, its findings are identical to those of Experiment 1. The second experiment (`Experiment 2b') was supposed to test the same hypothesis as Experiment 2, but the treatment--a very short distraction task--did not alter any of the outcome variables, suggesting that our outcome measures, a set of multiple-choice questions, was not sufficiently sensitive.

\subsection{Experiment 1}\label{experiment-1}

The main goal of Experiment 1 was to test whether exposure to impressive science content could enhance people's trust in scientists. We presented participants with vignettes about scientific findings in the disciplines of entomology and archaeology. We manipulated the impressiveness of the texts, so that there was one `basic' and one `impressive' version for each of the disciplines (see Table \ref{tab:exp1-stimuli}). We varied impressiveness within participants, but between disciplines: each participant was randomly assigned to see an impressive version for one discipline, and a basic version for the other discipline. We had the following hypotheses:

\textbf{H1a: Participants will perceive scientists as more competent than they did before, after having read an impressive text about their discipline's findings, compared to when reading a basic text.}

\textbf{H1b: Across all conditions, participants who are more impressed by the text about a discipline will also tend to perceive the scientists of the discipline as more competent.}

\textbf{H2a: Participants will trust a discipline more than they did before, after reading an impressive text about the discipline's findings, compared to when reading a basic text.}

\textbf{H2b: Across all conditions, participants who are more impressed by the text about a discipline will also tend to trust the scientists of the discipline more.}

We had the following research questions:

\textbf{RQ1: Do participants perceive that they learn more from the texts in the impressive condition, compared to the basic condition?}

\textbf{RQ2: Do perceptions of consensus interact with the relationships proposed in the hypotheses, such that greater perceived consensus is associated with a more positive relationship between impressiveness and trust/competence?}

\subsection{Methods}\label{methods}

\subsubsection{Procedure}\label{procedure}

After providing their consent to participate in the study and passing an attention check (see ESM), participants read the following introduction: ``You're going to read two texts about scientific discoveries, in two fields: archeology (the study of human activity through the recovery and analysis of material culture), and entomology (the study of insects). We ask you to read the texts, and then answer a few questions about them.'' Next, participants read two vignettes (one basic and one impressive) about scientific findings in the disciplines of entomology and archaeology. After reading each vignette, participants were asked a series of questions. We asked participants how much they thought they learned from reading the text (``How much do you feel you've learnt about human history by reading this text?'' {[}1 - Nothing, 2 - A bit, 3 - Some, 4 - Quite a bit, 5 - A lot{]}) and how impressive they found it (``How impressive do you think the findings of the archaeologists described in the text are?'' {[}1 - Not very impressive, 2 - A bit impressive, 3 - Quite impressive, 4 - Very impressive, 5 - Extremely impressive{]}). We also asked them whether reading the text changed their impression of the competence of the scientists of the respective discipline (``Would you agree that reading this text has made you think of archaeologists as more competent than you thought before?'' {[}1 - Strongly disagree, 2 - Disagree, 3 - Neither agree nor disagree, 4 - Agree, 5 - Strongly agree{]}) and their trust in the respective discipline (``Having read this text, would you agree that you trust the discipline of archaeology more than you did before?'' {[}1 - Strongly disagree, 2 - Disagree, 3 - Neither agree nor disagree, 4 - Agree, 5 - Strongly agree{]}). Finally, we also asked about the perceived consensus (``To which extent do you think the findings from the short text you just read reflect a minority or a majority opinion among archaeologists?'' {[}1 - Small minority, 2 - Minority, 3 - About half, 4 - Majority, 5 - Large majority{]}).

\subsubsection{Participants}\label{participants}

A power simulation (see OSF) suggested that the minimum required sample size to detect a statistically significant effect for all hypotheses with a power of 0.9 is n = 100 participants. We therefore recruited 100 participants from the UK via Prolific. One participant failed the attention check, resulting in a final sample of 99 participants (49 female, 50 male; \(age_\text{mean}\): 42.07, \(age_\text{sd}\): 12.86, \(age_\text{median}\): 40).

\subsubsection{Materials}\label{materials}

Table \ref{tab:exp1-stimuli} shows the stimuli used in Experiment 1.

\begingroup\fontsize{8}{10}\selectfont

\begin{longtable}[t]{>{\raggedright\arraybackslash}p{5em}>{\raggedright\arraybackslash}p{25em}>{\raggedright\arraybackslash}p{25em}}
\caption{\label{tab:exp1-stimuli}Stimuli of Experiment 1}\\
\toprule
 & Impressive & Basic\\
\midrule
\endfirsthead
\caption[]{\label{tab:exp1-stimuli}Stimuli of Experiment 1 \textit{(continued)}}\\
\toprule
 & Impressive & Basic\\
\midrule
\endhead

\endfoot
\bottomrule
\endlastfoot
Archeology & Archaeologists, scientists who study human history and prehistory, are able to tell, from their bones, whether someone was male or female, how old they were, and whether they suffered from a range of diseases. Archaeologists can now tell at what age someone, dead for tens of thousands of years, stopped drinking their mother’s milk, from the composition of their teeth.
Archaeologists learn about the language that our ancestors or cousins might have had. For instance, the nerve that is used to control breathing is larger in humans than in apes, plausibly because we need more fine-grained control of our breathing in order to speak. As a result, the canal containing that nerve is larger in humans than in apes – and it is also enlarged in Neanderthals.
Archaeologists can also tell, from an analysis of the tools they made, that most Neanderthals were right-handed. It’s thought that handedness is related to the evolution of language, another piece of evidence suggesting that Neanderthals likely possessed a form of language. & Archaeology is the science that studies human history and prehistory based on the analysis of objects from the past such as human bones, engravings, constructions, and various objects, from nails to bits of pots. This task requires a great deal of carefulness, because objects from the past need to often be dug out from the ground and patiently cleaned, without destroying them in the process.

Archaeologists have been able to shed light on human history in all continents, from ancient Egypt to the Incas in Peru or the Khmers in Cambodia.

Archaeologists study the paintings made by our ancestors, such as those that can be found in Lascaux, a set of caves found in the south of France that have been decorated by people at least 30000 years ago.

Archaeologists have also found remains of our more distant ancestors, showing that our species is just one among several that appeared, and then either changed or went extinct, such as Neanderthals, Homo erectus, or Homo habilis.\\
Entomology & Entomologists are the scientists who study insects. Some of them have specialized in understanding how insects perceive the world around them, and they have uncovered remarkable abilities. 

Entomologists interested in how flies’ visual perception works have used special displays to present images for much less than the blink of an eye, electrodes to record how individual cells in the flies’ brain react, and ultra-precise electron microscopy to examine their eyes. Thanks to these techniques, they have shown that some flies can perceive images that are displayed for just three milliseconds (a thousandth of a second) – about ten times shorter than a single movie frame (of which there are 24 per second). 

Entomologists who study the hair of crickets have shown that these microscopic hairs, which can be found on antenna-like organs attached to the crickets’ rear, are maybe the most sensitive organs in the animal kingdom. The researchers used extremely precise techniques to measure how the hair reacts to stimuli, such as laser-Doppler velocimetry, a technique capable of detecting the most minute of movements. They were able to show that the hair could react to changes in the motion of the air that had less energy than one particle of light, a single photon. & Entomologists are scientists who investigate insects, typically having a background in biology. They study, for example, how a swarm of bees organizes, or how ants communicate with each other. 

They also study how different insects interact with each other and their environment, whether some species are in danger of going extinct, or whether others are invasive species that need to be controlled.

Sometimes entomologists study insects by observing them in the wild, sometimes they conduct controlled experiments in laboratories, to see for example how different environmental factors change the behavior of insects, or to track exactly the same insects over a longer period of time.

An entomologist often specializes in one type of insect in order to study it in depth. For example, an entomologist who specializes in ants is called a myrmecologist.\\*
\end{longtable}
\endgroup{}

\subsection{Results and discussion}\label{results-and-discussion}

As a manipulation check, we find that participants indeed perceived the impressive texts to be more impressive (mean = 4.01, sd = 0.87; \(\hat{b}\) = 0.81 {[}0.569, 1.048{]}, p \textless{} .001) than the basic texts (mean = 3.20, sd = 1.13).

Participants perceived scientists as more competent after having read an impressive text (\(\hat{b}_{\text{Competence}}\) = 0.50 {[}0.314, 0.676{]}, p \textless{} .001; mean = 3.76, sd = 0.90) than after having read a basic one (mean = 3.26, sd = 0.80). Pooled across all conditions, participants impressiveness ratings were positively associated with competence: When participants reported being more impressed, they evaluated scientists as more competent (\(\hat{b}\) = 0.41 {[}0.312, 0.502{]}, p \textless{} .001).

Participants also trusted a discipline more after having read an impressive text (\(\hat{b}_{\text{trust}}\) = 0.32 {[}0.17, 0.476{]}, p \textless{} .001; mean = 3.68, sd = 0.87) than after having read a basic one. Participants impressiveness ratings were positively associated with trust when pooling across all conditions (\(\hat{b}\) = 0.30 {[}0.206, 0.385{]}, p \textless{} .001).

Regarding RQ1, participants had the impression of having learned more after having read the impressive text (mean = 3.66, sd = 0.93; \(\hat{b}\) = 0.89 {[}0.67, 1.108{]}, p \textless{} .001), than after having read the basic one (mean = 2.77, sd = 0.96).

Regarding RQ2, we do not find evidence that perceived consensus modulates the effect of impressiveness hypothesized in H1a and H2a (see ESM). In the ESM (section \ref{exp1}), we provide regression tables and figures.

\section{Experiment 2}\label{experiment-2}

In Experiments 1, exposure to impressive scientific content increased trust in the relevant science, and the perceived competence of the relevant scientists. In Experiment 2, we wanted to test whether these perceptions were at least partly independent of being able to recall the specific content that had induced them. Experiment 2 consisted of a `main study' and a `validation study'.

In the main study, each participant was randomly assigned to read the impressive version of one of the vignettes from Experiment 1 (Table \ref{tab:exp1-stimuli}). As in Experiment 1, participants were asked about how impressed they were by the findings and how whether they changed their perception of the scientists' competence and their trust in the scientific discipline.

In line with the findings of Experiment 1, we expected participants to have increased perceptions of competence and trust in scientists after:

\textbf{H1a: Participants perceive scientists as more competent after having read an impressive text about their discipline's findings, compared to before.}

\textbf{H1b: Participants trust a discipline more after having read an impressive text about the discipline's findings, compared to before.}

An addition of Experiment 2 was that participants were given a recall task: They were asked to rewrite the vignette text they've read just before from memory. We predicted that participants would not be able to recall all the information presented in the short vignettes right after having read them (see methods section for details):

\textbf{H2: Participants are not able to recall all the information of the original texts.}

This hypothesis, however, only tested whether people forgot any of the content, potentially including non-impressive content. To address this issue, we asked participants to select elements of the vignettes they found impressive (see Table \ref{tab:knowledge-evaluation-grid}). We predicted that even for the subjective subset of information that a participant said they found impressive, they forgot at least some of the content:

\textbf{H3: Participants are not able to recall all the impressive information--as rated by themselves--contained in the original text.}

H3 establishes whether participants recall or not content they themselves judge as impressive. However, this judgement could be biased as it happens after the recall task. For a more robust assessment, we conducted a validation study.

The validation study was run two days after the main study. For this study, a new sample of participants was recruited. Participants were randomly assigned to one of two conditions: In the ``original text'' condition, participants were assigned to read one of the two impressive texts of the main study, again picked at random. In the ``recall text'' condition, participants were also randomly assigned to either the archeology or the entomology condition, but instead of the original vignette texts, they read one of the recall texts written by the participants' of the main study. We predicted that participants in the validation study would be less impressed by the open-ended answers produced by the participants of Experiment 2 compared to the original impressive vignette texts (see Table \ref{tab:exp1-stimuli}) and, accordingly, would have less positively altered perceptions of the scientists' competence and the trustworthiness of their discipline:

\textbf{H4a: The texts produced by participants of the experiment as a result of the recall task will be less impressive than the original texts, as rated by participants of the validation study.}

\textbf{H4b: Participants of the validation study perceive scientists as more competent after having read the original texts, compared to after having read the texts produced by participants of the experiment as a result of the recall task.}

\textbf{H4c: Participants of the validation study trust a discipline more after having read the original texts, compared to after having read the texts produced by participants of the experiment as a result of the recall task.}

\subsection{Methods}\label{methods-1}

\subsubsection{Procedure}\label{procedure-1}

\paragraph{Main study}\label{main-study}

In the main study, after having consented to take part in the study and passing an attention check (see ESM), participants read the impressive version of one of the vignettes from Experiment 1 (Table \ref{tab:exp1-stimuli}). Whether The text was about archaeology or entomology was randomized. After reading the text, as in Experiment 1, participants were asked about changes in their perception of the scientists' competence (``Would you agree that reading this text has made you think of archaeologists/entomologists as more competent than you thought before?'' {[}1 = Strongly disagree, 2 = Disagree, 3 = Neither agree nor disagree, 4 = Agree, 5 = Strongly agree{]}) and trust in the scientists' discipline (``Having read this text, would you agree that you trust the discipline of archaeology/entomology more than you did before?'' {[}1 - Strongly disagree, 2 - Disagree, 3 - Neither agree nor disagree, 4 - Agree, 5 - Strongly agree{]}). They were also be asked about the impressiveness of the text they read: ``How impressive did you find the findings of the entomologists described in the text you have read?'' (5-point scale: 1 = not impressive at all, 2 = not very impressive, 3 = neither impressive nor unimpressive, 4 = quite impressive, 5 = very impressive). The order of these questions was randomized. Next, as an open-ended question, participants were asked to recall as much information as they could of the texts they had just read. They were told that their texts would be read by future participants. To further motivate participants, they were also told that they would get a bonus for recalling (without external aids) accurate information. They were not told how much that bonus would be. We payed them 5p per point gained in the recall task (see methods for how these points were assigned). This way, participants could reach a maximum bonus of 0.8 pound for archaeology (0.05p x 2 points x 8 content elements) and 0.7 pound (0.05p x 2 points x 7 content elements) for entomology. After that, participants were presented with the evaluation grid that we used to asses the open answers from the recall task (see Table \ref{tab:knowledge-evaluation-grid}). For each information element, we asked participants to indicate whether they found it impressive or not (``Do you think this piece of information is impressive?'' {[}Yes; No{]}). At the end of the main study, participants were asked about their education level.

\paragraph{Validation study}\label{validation-study}

After consenting to taking part in the study and passing an attention check (see ESM), participants read either one of the original vignettes, or a text produced as part of the recall task from the main study. For those participants assigned to read a recall text, the text was randomly sampled (with replacement) from a pool of recall answers. Orthographic and grammatical mistakes in these texts were corrected with the help of ChatGPT beforehand. We had preregistered to sample among all recall answers from participants of the main study. However, checking recall scores after the main study and before launching the validation study, we found them to be very low on average (see ESM section \ref{above-median}). To avoid having the answers of ``lazy'' participants--who might have recalled more, but were not motivated to type what they remembered--in our sample, we decided to only selected answers that scored above median in our recall measure for the validation study. This selection is conservative in that it makes it harder to confirm our predictions under H4.

After reading the text, just as in the main study, participants were asked about changes in their perception of the scientists' competence (e.g.~``Would you agree that reading this text has made you think of archaeologists/entomologists as more competent than you thought before?'' {[}1 = Strongly disagree, 2 = Disagree, 3 = Neither agree nor disagree, 4 = Agree, 5 = Strongly agree{]}) and trust in the scientists' discipline (``Having read this text, would you agree that you trust the discipline of archaeology/entomology more than you did before?'' {[}1 - Strongly disagree, 2 - Disagree, 3 - Neither agree nor disagree, 4 - Agree, 5 - Strongly agree{]}). Participants were also be asked about the impressiveness of the text they read: ``How impressive did you find the findings of the entomologists described in the text you have read?'' (5-point scale: 1 = not impressive at all, 2 = not very impressive, 3 = neither impressive nor unimpressive, 4 = quite impressive, 5 = very impressive). The order of these questions was randomized. Finally, participants were asked about their education level.

\subsubsection{Participants}\label{participants-1}

In our power simulation (see OSF), we varied the sample size of the main study and effect sizes (assuming the same effect size for all hypotheses in each scenario). For all simulations, we assumed a constant validation study sample size of 400 participants (only relevant for H4a, b and c). The power simulation suggested that we would reach a power level of 90\% when assuming a medium effect size of 0.5 with already 50 participants. Due to uncertainty about our assumptions, we recruited a sample of 203 participants for the main study, and a sample of 406 participants for the validation study. Although this was not the focus of our simulation, our results showed that for a medium effect size of 0.5, a sample size of 400 for the validation study yielded statistical power of greater than 90\% for all hypotheses based on this sample (H4a, b and c). All participants were from the UK, recruited via Prolific, and paid to complete the experiment.

The final sample comprised 198 participants (four failed attention checks; 99 female, 99 male; \(age_\text{mean}\): 42.65, \(age_\text{sd}\): 15.55, \(age_\text{median}\): 40) for the main study, 399 participants for the validation study (seven failed attention checks; 201 female, 198 male; \(age_\text{mean}\): 41.71, \(age_\text{sd}\): 13.48, \(age_\text{median}\): 40).

\subsubsection{Materials}\label{materials-1}

For Experiment 2, we used the impressive version of the stimuli used in Experiment 1 (see Table \ref{tab:exp1-stimuli}).

\begin{longtable}[t]{>{\raggedright\arraybackslash}p{20em}>{\raggedright\arraybackslash}p{20em}}
\caption{\label{tab:knowledge-evaluation-grid}Recall evaluation grid}\\
\toprule
Archaeology & Entomology\\
\midrule
1. Archaeologists can determine whether someone was male or female from their bones. & 1. Entomologists use special displays to present images to flies for extremely short periods (less than the blink of an eye).\\
2. Archaeologists can determine how old someone was from their bones. & 2. Entomologists can record how individual cells in flies’ brains react using electrodes.\\
3. Archaeologists can determine whether someone suffered from a range of diseases from their bones. & 3. Entomologists use ultra-precise electron microscopy to examine flies’ eyes.\\
4. Archaeologists can determine at what age someone stopped drinking their mother’s milk, based on the composition of their teeth. & 4. Some flies can perceive images displayed for just three milliseconds. This duration is about ten times shorter than a single movie frame.\\
5. The nerve controlling breathing is larger in humans than in apes. The canal containing that nerve is also larger in humans and Neanderthals than in apes. & 5. Crickets have microscopic hairs situated on antenna-like organs at their rear.\\
\addlinespace
6. The fact that the nerve controlling breathing is larger in humans is possibly due to the need for fine-grained control of breathing to speak. & 6. Crickets' hairs are possibly the most sensitive organs in the animal kingdom. They react to changes in air motion with less energy than one photon.\\
7. Archaeologists determined that most Neanderthals were right-handed, based on analysis of Neanderthals’ tools. & 7. Entomologists measured how cricket hairs react to stimuli, using laser-Doppler velocimetry, which can detect extremely minute movements.\\
8. Handedness is thought to be related to the evolution of language. This suggests that Neanderthals likely possessed a form of language. & \\
\bottomrule
\multicolumn{2}{l}{\rule{0pt}{1em}\textit{Note: }}\\
\multicolumn{2}{l}{\rule{0pt}{1em}For each knowledge element, participants could score a maximum of two points.}\\
\end{longtable}

\paragraph{Recall}\label{recall}

As shown in Table \ref{tab:knowledge-evaluation-grid}, we divided the text into a series of knowledge elements. For each participant, we coded to which extent they recalled each of the different elements. We then calculated a recall score based on how many of these elements they mentioned in their open-ended answer.

The coding was done with the help of ChatGPT (see ESM section \ref{gpt-prompt} for the exact prompt). The instructions were to code 0 if a piece of knowledge was not mentioned or was mentioned with very important mistakes (e.g., writing ``the nerve controlling fine hand movement is bigger in humans'' instead of ``the nerve controlling breathing is bigger in humans''). It should code 1 if the piece of knowledge was mentioned, but some important elements were missing (e.g., writing ``Archaeologists can determine at what age someone stopped drinking their mother's milk'' instead of « Archaeologists can determine at what age someone stopped drinking their mother's milk from the composition of their teeth »), and/or there were some mistakes (e.g., writing ``Archaeologists can determine at what age someone stopped drinking their mother's milk, based on the bones'' instead of ``Archaeologists can determine at what age someone stopped drinking their mother's milk based on the teeth''). It should code 2 if the piece of knowledge was mentioned with all the main content, even if the participant had not used the precise technical words (e.g., ``neanderthals'', ``laser-Doppler velocimetry'') or had changed the phrasing of the information in other ways. These instructions were intended to produce fairly generous recall scores.

Since the two vignettes contained a slightly different number of total information elements according to our evaluation grid (8 for archaeology, 7 for entomology), we used a relative measure for the final recall score, namely the share of obtained points among all possible points (possible range from 0 to 1). Practically, for all hypothesis tests on recall, we compute a forgetting score (1-recall score), and test whether this forgetting score is significantly different from zero.

To validate the scores assigned by ChatGPT, we compared them to scores assigned by two human coders for a subsample of 80 randomly chosen texts (half on archaeology, half on entomology). The human coders were unaware of the study context and the hypotheses. For instructions, they read the same prompt we provided to ChatGPT. To measure the reliability between ChatGPT and the human coders, we calculated an intraclass correlation coefficient (ICC) for two-way designs with random raters (Ten Hove, Jorgensen, \& Ark, 2024). Following the guidelines in Heyman, Lorber, Eddy, and West (2014), we preregistered taking 0.7 as a threshold for acceptable reliability. In our sample, we observe an ICC of 0.84 {[}se = 0.03{]}. To provide an intuition, human coders and ChatGPT, on average, have assigned exactly the same score in 72.9\% of all rating instances (compared to 73.5\% of exact agreement between the two human coders). We consider this a validation of using ChatGPT.

\paragraph{Recall of impressive items}\label{recall-of-impressive-items}

After giving their post-evaluation of scientists' competence and trust in the different disciplines, participants were presented with the evaluation grid shown in Table \ref{tab:knowledge-evaluation-grid} for the respective discipline they had been randomized to see. For each element in the evaluation grid, they were asked whether they found it impressive or not. For each participant, we computed the recall score just as described above, but on only on the subset of those elements they subjectively rated as impressive.

\subsection{Results and discussion}\label{results-and-discussion-1}

\subsubsection{Main study}\label{main-study-1}

First, as a kind of validation check, we note that most participants declared finding all knowledge elements to be impressive (\(mean_{\text{Archeology}}\) = 7.04, \(median_{\text{Archeology}}\) = 8, number of knowledge elements = 8; \(mean_{\text{Entomology}}\) = 6.39, \(median_{\text{Entomology}}\) = 7, number of knowledge elements = 7).

For the first set of hypotheses, we tested whether the average scores for perceived change in competence and trust were significantly different from their respective scale midpoint (which corresponds to no change in perception). We first tested the outcome variables' distributions for normality, using a Shapiro-Wilk test. In all cases, this test suggested that the data is considered non-normally distributed (p \textless{} 0.05). Following our pre-registration, we therefore did not run a default one-sample t-test, but used a Wilcoxon signed-rank test instead. We find evidence for both H1a and H1b: After having read an impressive text about the findings of a scientific discipline, participants saw the scientists as more competent (median = 4, W = 12,388.50, p \textless{} .001).

Despite having been impressed, a first descriptive analysis suggested that participants seemed to recall only very little information (\(mean\) = 32\%, \(sd\) = 23\%, \(median\) = 29\%). Given these low average scores, we opted for a more conservative approach to testing H2 and H3: We defined the median as a cut-off and selected only the 50\% of participants who remembered best. Our aim was to remove potentially ``lazy'' participants who did not make an effort. Since we hypothesized that participants would forget content, this selection made it harder for us to confirm our hypotheses. Even for the 50\% best participants, we find evidence for both hypotheses: Participants did not perfectly recall all information (H2, median = 0.43, W = 5778, p \textless{} .001) and they did not recall all information contained in the elements they judged as impressive themselves (H3, median = 0.43, W = 5778, p \textless{} .001)\footnote{The two median values for all information and for the subset of impressive information are not not exactly the same, but very similar, because participants rated most knowledge elements as impressive, making the general recall score and the recall score for impressive elements very similar.}.

\subsubsection{Validation study}\label{validation-study-1}

For H4a, b and c, we compared trust, competence and impressiveness ratings between the `original texts' and the `recall texts' conditions of the validation study using independent sample t-tests. Participants who read the original texts reported being more impressed (\(\hat{b}_{\text{Impressiveness}}\) = 0.39, t = 5.41, p \textless{} .001), rated the scientists of the respective discipline as more competent (\(\hat{b}_{\text{Competence}}\) = 0.27, t = 3.07, p = 0.002) and had more trust in the respective discipline (\(\hat{b}_{\text{Trust}}\) = 0.27, t = 3.40, p \textless{} .001) than participants who read the texts of the recall task.

\section{Discussion}\label{discussion}

It has long been a puzzle to the deficit model--which suggests that trust in science is primarily driven by science knowledge--that across many studies, knowledge about science is at best weakly associated with science attitudes. The rational impression account makes sense of this: It predicts that people trust science primarily because they have been impressed by it, not because they remember much of the knowledge that impressed them.

In two experiments, we have provided evidence for this account. Experiment 1 shows that impressive scientific findings lead people to think of scientists as more competent and trust science more. Experiment 2 shows that these impressions are formed even though participants forget the content that generated them almost immediately after reading it.

How can this account be rational, if it posits that trust is largely detached from knowledge? It is rational in that the heuristics it builds on lead to sound inferences in many contexts. If someone finds out something that is hard to know (and consensual, but here we take that for granted), such as the size of the Milky Way, we should expect them to be competent, even without knowing the details. Even forgetting specific knowledge is not irrational: It has been argued that one of the main functions of episodic memory is to justify our beliefs in communication with others (Mahr \& Csibra, 2018). As a consequence, we should be particularly good at remembering things we might need to convince others. In this regard, incentives of remembering science seem to be weak: Most exposure to science happens at school, and there is little reason for young learners' minds to anticipate having to convince others of the merits of science, when at school everyone readily accepts it.

Does this mean that scientific education should focus on impressive but potentially inaccessible findings, instead of fostering a proper understanding of more basic discoveries? No, for at least two reasons. First, for some students at least, this knowledge will be remembered, and will prove important in their lives. Second, positive impressions of science can also be created by a proper understanding of its explanations, methods, etc. -- even if, once again, that understanding is then often lost. However, our findings do suggest that science educators should not despair when they observe the low rates of science literacy in adults, since the exposure to science they provided is arguably responsible for creating a strong bedrock of trust in science.

The present framework also stresses the vital role that can be played by science communication. The relationship between science communication and trust in science has already been explored in depth König et al. (2023), but we believe the present framework might make a useful contribution. In particular, it suggests that, even though more understanding of the underlying methods is always preferable (König et al., 2023), even a relatively superficial exposure to impressive findings can bolster trust in science. An important caveat is that impressive findings that haven't yet gained the approval of the community might be particularly likely to backfire if people learn they have been disproven (on the importance of presenting a measure of consensus alongside scientific information, see König et al., 2023) .

The present experiments have a number of limitations: First, they were conducted on convenience samples recruited in a single country, the UK. Second, they were conducted within a very short time frame. While we can show that participants almost immediately forget about impressive content, it is not clear from our study for how long the impressions persist (although in other contexts impressions formed on the basis of a much more superficial exposure have been shown to last for months, Gunaydin, Selcuk, \& Zayas, 2017). Future studies could extend our findings to other populations and to longer time frames.

\subsubsection{Data availability}\label{data-availability}

Data for all experiments and the simulations is available on the OSF project page (). Note that on the project page, all materials related to what is referred to as ``Experiment 1'' in this paper are stored under ``experiment\_2'', and all materials related to ``Experiment 2'' in this paper are stored under ``experiment\_4''. This numbering is due to the original order in which experiments for this project were run. For a detailed report on the other experiments conducted as part of this project, please see the ESM.

\subsubsection{Code availability}\label{code-availability}

The code used to create all results (including tables and figures) of this manuscript is also available on the OSF project page ().

\subsubsection{Competing interest}\label{competing-interest}

The authors declare having no competing interests.

\FloatBarrier

\section{References}\label{references}

\phantomsection\label{refs}
\begin{CSLReferences}{1}{0}
\bibitem[\citeproctext]{ref-alganTrustScientistsTimes2021}
Algan, Y., Cohen, D., Davoine, E., Foucault, M., \& Stantcheva, S. (2021). Trust in scientists in times of pandemic: Panel evidence from 12 countries. \emph{Proceedings of the National Academy of Sciences}, \emph{118}(40), e2108576118. \url{https://doi.org/10.1073/pnas.2108576118}

\bibitem[\citeproctext]{ref-allumScienceKnowledgeAttitudes2008}
Allum, N., Sturgis, P., Tabourazi, D., \& Brunton-Smith, I. (2008). Science knowledge and attitudes across cultures: a meta-analysis. \emph{Public Understanding of Science}, \emph{17}(1), 35--54. \url{https://doi.org/10.1177/0963662506070159}

\bibitem[\citeproctext]{ref-bakEducationPublicAttitudes2001}
Bak, H.-J. (2001). Education and Public Attitudes toward Science: Implications for the {``}Deficit Model{''} of Education and Support for Science and Technology. \emph{Social Science Quarterly}, \emph{82}(4), 779--795. \url{https://doi.org/10.1111/0038-4941.00059}

\bibitem[\citeproctext]{ref-bauerWhatCanWe2007}
Bauer, M. W., Allum, N., \& Miller, S. (2007). What can we learn from 25 years of PUS survey research? Liberating and expanding the agenda. \emph{Public Understanding of Science}, \emph{16}(1), 79--95. \url{https://doi.org/10.1177/0963662506071287}

\bibitem[\citeproctext]{ref-beckRiskSocietyNew1992}
Beck, U. (1992). \emph{Risk society: towards a new modernity} (Repr). London: Sage.

\bibitem[\citeproctext]{ref-colognaRoleTrustClimate2020}
Cologna, V., \& Siegrist, M. (2020). The role of trust for climate change mitigation and adaptation behaviour: A meta-analysis. \emph{Journal of Environmental Psychology}, \emph{69}, 101428. \url{https://doi.org/10.1016/j.jenvp.2020.101428}

\bibitem[\citeproctext]{ref-durantPublicUnderstandingScience1989}
Durant, J. R., Evans, G. A., \& Thomas, G. P. (1989). The public understanding of science. \emph{Nature}, \emph{340}(6228), 11--14. \url{https://doi.org/10.1038/340011a0}

\bibitem[\citeproctext]{ref-feinsteinSustainedExperienceEmotion2010}
Feinstein, J. S., Duff, M. C., \& Tranel, D. (2010). Sustained experience of emotion after loss of memory in patients with amnesia. \emph{Proceedings of the National Academy of Sciences}, \emph{107}(17), 7674--7679. \url{https://doi.org/10.1073/pnas.0914054107}

\bibitem[\citeproctext]{ref-funkPublicConfidenceScientists2020}
Funk, C., \& Kennedy, B. (2020). \emph{Public confidence in scientists has remained stable for decades}.

\bibitem[\citeproctext]{ref-funkScienceScientistsHeld2020}
Funk, C., Tyson, A., Kennedy, B., \& Johnson, C. (2020). \emph{Science and scientists held in high esteem across global publics}. Retrieved from \url{https://www.pewresearch.org/science/2020/09/29/science-and-scientists-held-in-high-esteem-across-global-publics/}

\bibitem[\citeproctext]{ref-gauchatCulturalAuthorityScience2011}
Gauchat, G. (2011). The cultural authority of science: Public trust and acceptance of organized science. \emph{Public Understanding of Science}, \emph{20}(6), 751--770. \url{https://doi.org/10.1177/0963662510365246}

\bibitem[\citeproctext]{ref-gauchatCulturalCognitiveMappingScientific2018}
Gauchat, G., \& Andrews, K. T. (2018). The Cultural-Cognitive Mapping of Scientific Professions. \emph{American Sociological Review}, \emph{83}(3), 567--595. \url{https://doi.org/10.1177/0003122418773353}

\bibitem[\citeproctext]{ref-giddensModernitySelfidentitySelf1991}
Giddens, A. (1991). \emph{Modernity and self-identity: Self and society in the late modern age}. Stanford, Calif: Stanford University Press.

\bibitem[\citeproctext]{ref-gunaydinImpressionsBasedPortrait2017}
Gunaydin, G., Selcuk, E., \& Zayas, V. (2017). Impressions Based on a Portrait Predict, 1-Month Later, Impressions Following a Live Interaction. \emph{Social Psychological and Personality Science}, \emph{8}(1), 36--44. \url{https://doi.org/10.1177/1948550616662123}

\bibitem[\citeproctext]{ref-habermasJurgenHabermasSociety1989}
Habermas, J. (1989). \emph{Jürgen Habermas on society and politics: a reader} (1. {[}print.{]}; S. Seidman, Ed.). Boston: Beacon Press.

\bibitem[\citeproctext]{ref-heymanBehavioralObservationCoding2014}
Heyman, R. E., Lorber, M. F., Eddy, J. M., \& West, T. V. (2014). \emph{Behavioral observation and coding} (C. M. Judd \& H. T. Reis, Eds.). Cambridge: Cambridge University Press. \url{https://doi.org/10.1017/CBO9780511996481.018}

\bibitem[\citeproctext]{ref-konigHowCommunicateScience2023}
König, L. M., Altenmüller, M. S., Fick, J., Crusius, J., Genschow, O., \& Sauerland, M. (2023). \emph{How to communicate science to the public? Recommendations for effective written communication derived from a systematic review}. \url{https://doi.org/10.31234/osf.io/cwbrs}

\bibitem[\citeproctext]{ref-lindholtPublicAcceptanceCOVID192021}
Lindholt, M. F., Jørgensen, F., Bor, A., \& Petersen, M. B. (2021). Public acceptance of COVID-19 vaccines: cross-national evidence on levels and individual-level predictors using observational data. \emph{BMJ Open}, \emph{11}(6), e048172. \url{https://doi.org/10.1136/bmjopen-2020-048172}

\bibitem[\citeproctext]{ref-liquinMotivatedLearnAccount2022}
Liquin, E. G., \& Lombrozo, T. (2022). Motivated to learn: An account of explanatory satisfaction. \emph{Cognitive Psychology}, \emph{132}, 101453. \url{https://doi.org/10.1016/j.cogpsych.2021.101453}

\bibitem[\citeproctext]{ref-lombrozoStructureFunctionExplanations2006}
Lombrozo, Tania. (2006). The structure and function of explanations. \emph{Trends in Cognitive Sciences}, \emph{10}(10), 464--470. \url{https://doi.org/10.1016/j.tics.2006.08.004}

\bibitem[\citeproctext]{ref-lombrozoSimplicityProbabilityCausal2007}
Lombrozo, T. (2007). Simplicity and probability in causal explanation. \emph{Cognitive Psychology}, \emph{55}(3), 232--257. \url{https://doi.org/10.1016/j.cogpsych.2006.09.006}

\bibitem[\citeproctext]{ref-mahrWhyWeRemember2018}
Mahr, J. B., \& Csibra, G. (2018). Why do we remember? The communicative function of episodic memory. \emph{Behavioral and Brain Sciences}, \emph{41}, e1. \url{https://doi.org/10.1017/S0140525X17000012}

\bibitem[\citeproctext]{ref-millerMeasurementCivicScientific1998a}
Miller, J. D. (1998). The measurement of civic scientific literacy. \emph{Public Understanding of Science}, \emph{7}(3), 203--223. \url{https://doi.org/10.1088/0963-6625/7/3/001}

\bibitem[\citeproctext]{ref-millerPublicUnderstandingAttitudes2004}
Miller, J. D. (2004). Public Understanding of, and Attitudes toward, Scientific Research: What We Know and What We Need to Know. \emph{Public Understanding of Science}, \emph{13}(3), 273--294. \url{https://doi.org/10.1177/0963662504044908}

\bibitem[\citeproctext]{ref-nationalacademiesofsciencesengineeringandmedicineScienceLiteracyConcepts2016}
National Academies of Sciences, Engineering, and Medicine. (2016). \emph{Science Literacy: Concepts, Contexts, and Consequences} (C. E. Snow \& K. A. Dibner, Eds.). Washington, D.C.: National Academies Press. \url{https://doi.org/10.17226/23595}

\bibitem[\citeproctext]{ref-noyScienceGoodEffects2019}
Noy, S., \& O'Brien, T. L. (2019). Science for good? The effects of education and national context on perceptions of science. \emph{Public Understanding of Science}, \emph{28}(8), 897--916. \url{https://doi.org/10.1177/0963662519863575}

\bibitem[\citeproctext]{ref-pfanderHowWiseCrowd2025}
Pfänder, J., De Courson, B., \& Mercier, H. (2025). How wise is the crowd: Can we infer people are accurate and competent merely because they agree with each other? \emph{Cognition}, \emph{255}, 106005. \url{https://doi.org/10.1016/j.cognition.2024.106005}

\bibitem[\citeproctext]{ref-pfanderQuasiuniversalAcceptanceBasic2024}
Pfänder, J., Kerzreho, L., \& Mercier, H. (2024). \emph{Quasi-universal acceptance of basic science in the US}. \url{https://doi.org/10.31219/osf.io/qc43v}

\bibitem[\citeproctext]{ref-readExplanatoryCoherenceSocial1993}
Read, S. J., \& Marcus-Newhall, A. (1993). Explanatory coherence in social explanations: A parallel distributed processing account. \emph{Journal of Personality and Social Psychology}, \emph{65}(3), 429--447. \url{https://doi.org/10.1037/0022-3514.65.3.429}

\bibitem[\citeproctext]{ref-smithTrendsPublicAttitudes2013}
Smith, T. W., \& Son, J. (2013). Trends in public attitudes about confidence in institutions. \emph{General Social Survey Final Report}.

\bibitem[\citeproctext]{ref-sturgisTrustScienceSocial2021}
Sturgis, P., Brunton-Smith, I., \& Jackson, J. (2021). Trust in science, social consensus and vaccine confidence. \emph{Nature Human Behaviour}, \emph{5}(11), 1528--1534. \url{https://doi.org/10.1038/s41562-021-01115-7}

\bibitem[\citeproctext]{ref-tenhoveUpdatedGuidelinesSelecting2024}
Ten Hove, D., Jorgensen, T. D., \& Ark, L. A. van der. (2024). Updated guidelines on selecting an intraclass correlation coefficient for interrater reliability, with applications to incomplete observational designs. \emph{Psychological Methods}, \emph{29}(5), 967--979. \url{https://doi.org/10.1037/met0000516}

\bibitem[\citeproctext]{ref-weingartScienceCommunicationIssue2016}
Weingart, P., \& Guenther, L. (2016). Science communication and the issue of trust. \emph{Journal of Science Communication}, \emph{15}(05), C01. \url{https://doi.org/10.22323/2.15050301}

\bibitem[\citeproctext]{ref-wellcomeglobalmonitorWellcomeGlobalMonitor2018}
Wellcome Global Monitor. (2018). \emph{Wellcome Global Monitor 2018}. Retrieved from \url{https://wellcome.org/reports/wellcome-global-monitor/2018}

\bibitem[\citeproctext]{ref-wellcomeglobalmonitorWellcomeGlobalMonitor2020}
Wellcome Global Monitor. (2020). \emph{Wellcome Global Monitor 2020: Covid-19}. Retrieved from \url{https://wellcome.org/reports/wellcome-global-monitor-covid-19/2020}

\end{CSLReferences}

\newpage

\appendix



\end{document}
